\documentclass[12pt]{article}
\usepackage[T1]{fontenc}
\usepackage[utf8]{inputenc}
\usepackage{lmodern}
\usepackage{textcomp}
\usepackage{enumitem}
\usepackage{geometry}
\geometry{margin=1in}
\begin{document}
\begin{center}
Answer Key - Political Science - Undergraduate Level Assessment 17 \
\small (Answers are ordered exactly as in the question paper.)
\end{center}

\section*{Multiple Choice}
\begin{enumerate}[label=\arabic*.]
\item \textbf{Answer:} B
\\ \textit{Options:} A) It would have drained the US economically; B) It would have involved compromising opposition to slavery; C) It would have increased immigration to the US; D) None of the above
\\ \textit{Note:} Lincoln opposed the southward expansion of the US because such expansion would necessitate a compromise on the issue of slavery, which he vehemently opposed as it would undermine the principles of freedom and equality.
\item \textbf{Answer:} D
\\ \textit{Options:} A) Appellate judges never use the stare decisis principle.; B) Appellate trials are usually jury trials.; C) Appellate judges usually decide the facts of a case.; D) Appellate judges review decisions of lower courts.
\\ \textit{Note:} Appellate judges are tasked with reviewing and evaluating the decisions made by lower courts to ensure that the law was applied correctly and to address any legal errors that may have occurred during the initial trial.
\item \textbf{Answer:} D
\\ \textit{Options:} A) Any crime that uses computers to jeopardise or attempt to jeopardise national security; B) The use of computer networks to commit financial or identity fraud; C) The theft of digital information; D) Any crime that involves computers and networks
\\ \textit{Note:} The term 'cyber-crime' refers to illegal activities that are conducted through the use of computers and networks, encompassing a wide range of offenses such as hacking, identity theft, and online fraud.
\item \textbf{Answer:} D
\\ \textit{Options:} A) Framers' unqualified commitment to individual rights; B) small states' determination to receive equal representation in the legislature; C) Northern states' support for abolitionism; D) states' fears of an overpowerful national government
\\ \textit{Note:} The swift adoption of the Bill of Rights reflects the states' concerns about the potential for federal overreach and the need to protect individual liberties against an overpowerful national government, which was a central issue during the ratification debates of the Constitution.
\item \textbf{Answer:} D
\\ \textit{Options:} A) Its primary role is to appropriate spending for infrastructure projects.; B) It is the primary author of congressional banking reform legislation.; C) It allocates funding for canals and waterways.; D) Its jurisdiction includes the tax system.
\\ \textit{Note:} The House Committee on Ways and Means is responsible for overseeing all taxation and revenue-related legislation, making it the primary committee in Congress that addresses issues related to the tax system.
\end{enumerate}

\section*{True / False}
\begin{enumerate}[label=\arabic*.]
\setcounter{enumi}{5}
\item \textbf{Answer:} True
\\ \textit{Options:} A) True; B) False
\\ \textit{Note:} Divided government often results in political gridlock, where differing priorities and party affiliations between the executive and legislative branches can hinder the timely confirmation of ambassadorial nominees.
\item \textbf{Answer:} True
\\ \textit{Options:} A) True; B) False
\\ \textit{Note:} The Strategic Arms Reduction Treaty (START) was designed to reduce and limit the number of strategic offensive arms, specifically long-range nuclear missiles, thereby mandating their elimination between the United States and the Soviet Union.
\item \textbf{Answer:} False
\\ \textit{Options:} A) True; B) False
\\ \textit{Note:} Delegates tend to be more aligned with the major parties, reflecting the preferences of their party's base, which makes them less likely to register as third party voters compared to the general population.
\item \textbf{Answer:} False
\\ \textit{Options:} A) True; B) False
\\ \textit{Note:} The legislative veto was a mechanism that allowed Congress to override executive actions, not for the executive branch to veto legislation passed by Congress.
\item \textbf{Answer:} True
\\ \textit{Options:} A) True; B) False
\\ \textit{Note:} The statement is true because a child's family serves as their primary socialization agent, influencing their beliefs, attitudes, and political orientations through direct communication, shared experiences, and the transmission of cultural norms and values.
\end{enumerate}

\section*{Long Form Response}
\begin{enumerate}[label=\arabic*.]
\setcounter{enumi}{10}
\item \textbf{Answer:} The Ideal Policy framework serves as a valuable analytical tool in political science by focusing on coercive diplomacy as a means to counter aggression. This framework emphasizes the importance of the coercer's actions and outlines how these actions can lead to successful outcomes with a minimal set of success conditions. By predicting the effectiveness of coercive strategies, the Ideal Policy framework allows practitioners to assess various diplomatic scenarios and their potential impacts. This predictive capability is essential for crafting effective foreign policy responses, as it enables policymakers to anticipate the reactions of adversaries and make informed decisions based on the likelihood of achieving desired objectives. Overall, the Ideal Policy framework enhances our understanding of coercive diplomacy and equips practitioners with the insights necessary for effective engagement in international relations.
\\ \textit{Note:} correct because it highlights the Ideal Policy framework's focus on coercive diplomacy and its predictive capabilities, emphasizing how it enables practitioners to evaluate the effectiveness of diplomatic strategies in countering aggression based on clear success conditions.
\item \textbf{Answer:} Postcolonialism and security studies intersect in a complex and often contentious manner, as postcolonial critiques challenge the traditional Westphalian perspective that underpins much of security analysis. The Westphalian model emphasizes state sovereignty and territorial integrity, which can overlook the historical and ongoing impacts of colonialism and imperialism on global security dynamics. Postcolonial theorists argue that this perspective is limited because it fails to account for the experiences and voices of formerly colonized nations, which often face different security threats rooted in historical injustices. By highlighting the inadequacies of the Westphalian framework, postcolonialism calls for a re-examination of security that includes cultural, social, and economic dimensions, thereby expanding the scope of security studies to consider a more inclusive and comprehensive understanding of global relations.
\\ \textit{Note:} correct because it identifies how postcolonial critiques of the Westphalian perspective reveal its inadequacies in addressing the historical and contemporary security issues rooted in colonialism and imperialism, thereby enriching the analysis of global security dynamics.
\end{enumerate}

\vfill
\noindent\textit{End of Answer Key}
\end{document}